\chapter{Theory}

The first paragraph of the theory chapter in your Honours thesis plays a crucial role in setting the stage for the theoretical framework you're presenting, but it’s important not to rush it. It is best to write this part last, after you have completed all the sections and subsections of your theory chapter. This allows you to have a clear understanding of the key themes, concepts, and arguments you’ve covered, ensuring that your introduction to the theory chapter accurately reflects and ties together all the information. In this opening paragraph, you should briefly introduce the core theoretical concepts that will be explored and provide an overview of how these theories connect to your research question. By waiting until your theory chapter is fully developed, you’ll be able to craft a more cohesive and compelling introduction that logically leads the reader into the complex discussions that follow.

\section{Name of First Section}
Lorem ipsum dolor sit amet, consectetur adipiscing elit. Vivamus in tincidunt libero. Nulla facilisi. Nunc sit amet velit nec turpis suscipit pharetra. Sed convallis turpis vitae quam luctus, nec varius arcu consequat. Suspendisse potenti. Nulla euismod semper libero, a consequat magna eleifend vel. Curabitur lacinia dolor a ante pharetra, in tincidunt purus volutpat. Perseverare in omnibus, nam fortitudo et virtus in patientia consistunt. Mauris ut sem libero. Aliquam erat volutpat. Proin sollicitudin purus ut auctor feugiat.

\subsection{Name of First Sections First Subsection}
Knock, knock.
Who's there?
Cite.
Cite who?
\subsection{Name of First Sections Second Subsection}
Cite your sources, or your research is incomplete!
\section{Name of Second Section}

\subsection{Some Physics Equations}

Newton's Second Law of Motion:
\begin{equation}
F = ma
\end{equation}

Einstein's Mass-Energy Equivalence:
\begin{equation}
E=mc^2
\end{equation}

The Schrödinger Equation:
\begin{equation}
i\hbar \frac{\partial \Psi}{\partial t} = -\frac{\hbar^2}{2m} \nabla^2 \Psi + V \Psi
\end{equation}

Maxwell's Equations:
\begin{align}
\nabla \cdot \mathbf{E} &= \frac{\rho}{\epsilon_0} \\
\nabla \cdot \mathbf{B} &= 0 \\
\nabla \times \mathbf{E} &= -\frac{\partial \mathbf{B}}{\partial t} \\
\nabla \times \mathbf{B} &= \mu_0\mathbf{J} + \mu_0\epsilon_0 \frac{\partial \mathbf{E}}{\partial t}
\end{align}

\subsection{Some Mathematics Equations}

The Quadratic Formula:
\begin{equation}
x = \frac{-b \pm \sqrt{b^2 - 4ac}}{2a}
\end{equation}

The Fundamental Theorem of Calculus:
\begin{equation}
\int_a^b f(x) \, dx = F(b) - F(a)
\end{equation}

Euler's Identity:
\begin{equation}
e^{i\pi} + 1 = 0
\end{equation}

The Pythagorean Theorem:
\begin{equation}
a^2 + b^2 = c^2
\end{equation}

\section{Name of Third Section}
\subsection{Some Statistics Equations}

The Mean of a Data Set:
\begin{equation}
\bar{x} = \frac{1}{n} \sum_{i=1}^{n} x_i
\end{equation}

The Variance of a Data Set:
\begin{equation}
\sigma^2 = \frac{1}{n-1} \sum_{i=1}^{n} (x_i - \bar{x})^2
\end{equation}

The Normal Distribution Probability Density Function:
\begin{equation}
f(x | \mu, \sigma^2) = \frac{1}{\sqrt{2\pi\sigma^2}} e^{-\frac{(x-\mu)^2}{2\sigma^2}}
\end{equation}

\subsection{Some Engineering Equations}

Ohm's Law for Electrical Circuits:
\begin{equation}
V = IR
\end{equation}

Hooke's Law for Elasticity:
\begin{equation}
F = k\Delta L
\end{equation}

The Navier-Stokes Equation for Fluid Dynamics:
\begin{equation}
\rho \left( \frac{\partial \mathbf{v}}{\partial t} + \mathbf{v} \cdot \nabla \mathbf{v} \right) = -\nabla P + \mu \nabla^2 \mathbf{v} + \mathbf{f}
\end{equation}

\subsection{Some Computer Science Equations}

The Time Complexity of Bubble Sort:
\begin{equation}
O(n^2)
\end{equation}

The Binary Search Algorithm:
\begin{equation}
\text{BinarySearch}(A, x, \text{low}, \text{high}) = 
\begin{cases}
\text{True} & \text{if } A[\text{mid}] = x \\
\text{False} & \text{if } \text{low} > \text{high} \\
\text{BinarySearch}(A, x, \text{low}, \text{mid} - 1) & \text{if } A[\text{mid}] > x \\
\text{BinarySearch}(A, x, \text{mid} + 1, \text{high}) & \text{otherwise}
\end{cases}
\end{equation}

These equations encompass a wide range of mathematical and scientific concepts. You can use these LaTeX code snippets as a starting point and customize them to suit the specific equations you need for your research or document.